
The approach comprises three stages: (1) compute semantic entropy distributions per benchmark, (2) cluster benchmarks by JS-divergence, and (3) evaluate transfer success within and across clusters.

\subsubsection{3.1 Semantic Entropy Computation}

Following Kuhn et al. (2023), semantic entropy is computed for each query as follows:

\textbf{Generation.} For each query $q, N=10$ responses $\{r_1, \ldots, r_N\}$ are generated using temperature $T=0.7$ (or $T=1.0$ for transfer experiments) sampling from Llama-2-7B-Chat.

\textbf{Semantic clustering.} Responses are clustered by meaning using bidirectional entailment. Two responses $r_i, r_j$ belong to the same semantic cluster if $\text{NLI}(r_i, r_j) = \text{ENTAILMENT}$ and $\text{NLI}(r_j, r_i) = \text{ENTAILMENT}$, using DeBERTa-v3-large-MNLI as the NLI model.

\textbf{Entropy computation.} Let $C_1, \ldots, C_k$ be the semantic clusters with empirical probabilities $p_c = |C_c|/N$. Semantic entropy is:
$$H_{\text{sem}}(q) = -\sum_{c=1}^{k} p_c \log p_c$$

High entropy indicates diverse semantic content across generations; low entropy indicates consistent responses.

\subsubsection{3.2 Distribution Distance via JS-Divergence}

To quantify similarity between benchmark uncertainty distributions, Jensen-Shannon divergence is used. For benchmarks $B_i$ and $B_j$ with semantic entropy distributions $P_i$ and $P_j$:
$$\text{JS}(P_i \| P_j) = \frac{1}{2} D_{\text{KL}}(P_i \| M) + \frac{1}{2} D_{\text{KL}}(P_j \| M)$$
where $M = \frac{1}{2}(P_i + P_j)$.

Distributions are estimated using kernel density estimation (KDE) with Gaussian kernels. JS-divergence is symmetric, bounded in $[0, 1]$, and interpretable: values near 0 indicate similar distributions; values near 1 indicate dissimilar distributions.

\subsubsection{3.3 Hierarchical Clustering}

Benchmark families are discovered using hierarchical agglomerative clustering with Ward linkage on the JS-divergence matrix. This approach does not require pre-specifying the number of clusters and produces interpretable dendrograms. The optimal cluster count is selected by maximizing silhouette score.

\subsubsection{3.4 Transfer Evaluation Protocol}

\textbf{Threshold calibration.} For source benchmark $B_s$, data is split 70/30 into calibration and held-out sets. A threshold $\tau_s$ is calibrated to achieve 10\% false positive rate on the calibration set.

\textbf{Transfer evaluation.} Threshold $\tau_s$ is applied to target benchmark $B_t$ and AUROC is measured. Transfer degradation is:
$$\Delta_{\text{AUROC}} = \text{AUROC}(B_s) - \text{AUROC}(B_t | \tau_s)$$

The pre-specified criteria were: within-cluster degradation $\leq 0.08$ and cross-cluster degradation $> 0.15$.

